% !TeX root = ../教材题库_全部条目_转写工作稿.tex
% Source: XJ-B1-C03-EX-P084-Y0177
% 第3章 / 3.2 函数的基本性质 / 教材 p.84 / PDF 第91页
% OCR未校对：公式与文字均可能有误；未完成数学TeX转写；禁止直接发给学生
% 原题及相邻原解截图（核对用）：../题目原图/XJ-B1-C03-EX-P084-Y0177-01.webp
% 本文件只在转写完成后作为题干使用；OCR 条目尚不是可靠题干。
\begin{bankitem}{3.2 例题 例题4}{第3章 / 3.2 函数的基本性质 / 教材 p.84 / PDF 第91页}{XJ-B1-C03-EX-P084-Y0177}{OCR未校对：公式与文字均可能有误；未完成数学TeX转写；禁止直接发给学生}
\begin{ocrdraft}
[教材p.84 / PDF第91页]
\par
0互判断下列函数的奇偶性：
\par
（D丿、(x)一x2十》到；
\par
（2）g（工）一工十
\par
解
\par
（D函数丆0）的定义域为R.因为对VxeR,均有—xeR，且不一到一
\par
（一到2十》一到一十]到一丆0），所以不x）为偶函数．
\par
（2）函数g(x）的定义域为00学0\}．因为对Ve0\}，均有—xe0
\par
工乒0\}，且g（一工）一一工十
\par
一g(x），所以g(x）为奇函数．
\par
（3）因为方(x）的定义域关于原点不对称，所以h(x）既不是奇函数也不是偶函数．
\end{ocrdraft}
\end{bankitem}
